\section{Specifications and Projection Detail}\label{app:specifications}

\subsection{Employment specifications}

The first of the three exercises in \cref{sec:specs} expresses the compositional gap of \cref{eq:shiftshare} as an estimable gradient. I stack occupation $\times$ region cells and estimate
\begin{equation}\label{eq:employment}
y_{k,r} = \alpha + \beta \Ebar_k + \delta\, \Scot_r +
\zeta \bigl( \Ebar_k \times \Scot_r \bigr) + \phi_{s(k)} + u_{k,r},
\end{equation}
weighted by cell employment, where $y_{k,r}$ is log occupational employment and, in a second specification, the employment share; $\phi_{s(k)}$ are one-digit SOC major-group fixed effects; and errors are clustered by occupation. Since $\Ebar_k$ does not vary by region, the interaction $\zeta$ is identified purely from the interaction of the common exposure score with the regional employment distribution, and a negative estimate means composition insulates Scotland.

The estimates agree with the decomposition they restate. With major-group fixed effects, $\widehat{\zeta}$ is $-0.629$ (SE 0.241) for log employment and $-0.018$ (0.007) for the employment share, impliying within broad occupation groups, Scottish employment sits systematically further down the exposure gradient than rUK employment. 

\subsection{Productivity calibration}

\paragraph{Where \cref{eq:tfp} comes from.} The projection rests on an envelope result rather than on an assumed functional form, and deriving it matters here because it is what makes the regional differential a linear functional of the scores and therefore subject to \cref{cor:functionals}. Let the final good be produced under constant returns from a set of tasks with prices $p_1, \dots, p_N$ and unit cost function $c(p_1, \dots, p_N)$, homogeneous of degree one. By Shephard's lemma the cost-minimising task demands are $\partial c / \partial p_n$, so
\begin{equation}\label{eq:shares}
\varsigma_n \;=\; \frac{\partial \ln c}{\partial \ln p_n}
\;=\; \frac{p_n \, (\partial c / \partial p_n)}{c}
\end{equation}
is the Shephard cost share of task $n$, and the shares sum to one by Euler's theorem under linear homogeneity. If AI lowers the cost of task $n$ by the fraction $\kappa_n$, so that $d \ln p_n = -\kappa_n$, then totally differentiating gives $d \ln c = \sum_n \varsigma_n \, d\ln p_n = - \sum_n \varsigma_n \kappa_n$, and since total factor productivity is the inverse of unit cost at fixed input prices,
\begin{equation}\label{eq:hulten}
\Delta \ln \mathrm{TFP} \;=\; -\, d \ln c \;=\; \sum_n \varsigma_n \, \kappa_n .
\end{equation}
This is \citet{hulten1978} applied at the task rather than the sector level, which is the step \citet{acemoglu2025} takes within the task framework of \citet{acemoglu2019}. Aggregating \cref{eq:hulten} over a region's employment distribution requires one further step, and it is where the second parameter of \cref{eq:tfp} enters. Set $\kappa_n = \kappa$ on tasks AI in fact performs and zero otherwise. Exposure identifies the task content AI \emph{could} perform, and only a fraction $\pi$ of that content is profitably automated over the horizon, so the displaced cost share of minor group $k$ is $\pi \Ebar_k$ rather than $\Ebar_k$. Substituting into \cref{eq:hulten} and weighting by employment gives \cref{eq:tfp}. The two parameters are therefore distinct objects: $\varsigma_n$ in \cref{eq:shares} is a Shephard cost share that sums to one across tasks, whereas $\pi$ in \cref{eq:tfp} is an automatable fraction bounded in $[0,1]$ and constant across occupations. Only the latter is calibrated.

\paragraph{The imported parameters.} I import deep params, so that the Scottish and rUK paths differ only through their employment-weighted exposure. The exposure measurement is my own and the architecture follows the \citet{obr2025}, but the productivity parameters are Acemoglu's rather than the OBR's, with one exception noted below. The average cost saving on an exposed task, $\kappa$, rests on the two task-level randomised experiments Acemoglu relies on: \citet{noy2023}, who find substantial time savings on mid-level professional writing tasks, and \citet{brynjolfsson2025}, who find productivity gains among customer-support agents concentrated among the least experienced. \citet{noy2023} find a 40\% labour cost saving on writing tasks and \citet{brynjolfsson2025} find 14\% on customer-support tasks; averaging the two gives labour cost savings of about 27\% on the tasks they affect, which \citet{acemoglu2025} converts to overall cost savings of about 15\% using industry labour shares. The low, central and high columns of \cref{tab:projection} move $\kappa$ across $\{0.110,\, 0.144,\, 0.182\}$, and all three are overall rather than labour cost savings. The central setting is the value implied by Acemoglu's own arithmetic: $0.0066 / (0.199 \times 0.23) = 0.144$, so it reproduces his headline 0.66\% at his exposure share and automatable fraction. The low setting is the bottom of his 0.53--0.71\% range. The high setting is the \citet{obr2025} figure, which takes a task-level cost saving of 32\% and a UK labour share of 0.57, giving $0.32 \times 0.57 = 0.182$; the OBR pairs it with a more generous automatable fraction of 0.33 and an exposure share of 0.40, which I do not adopt, so only the cost-saving parameter is imported here. Applying Acemoglu's 27\% \emph{labour} cost saving directly as $\kappa$ would double-count the labour share and is not a scenario but a change of units. The automatable fraction is set to $\pi_k = \pi = 0.23$ for every occupation, this being the share of exposed tasks that Acemoglu judges profitably automatable over the decade. Because $\pi_k$ does not vary across $k$, \cref{cor:functionals} applies and the regional differential inherits the cancellation of \cref{prop:cancellation} exactly; the levels of the path, by contrast, are proportional to $\pi$ and carry its full model dependence. I retain the ten-year horizon. The contribution is the regional differential in the implied path, not a reassessment of Acemoglu's headline magnitude, which is contested and outside the scope of this dissertation. The reported regional levels (2.21\% for rUK, 2.18\% for Scotland) sit above his national figure of about 0.66\% by construction, because my continuous index is broader than his roughly 20\% task-exposure measure, so I place weight on the differential and its sign, and none on the levels.

\subsection{The elasticity of NSND revenue to earnings}\label{app:elasticity}

The third step of \cref{sec:fiscal} needs the proportional change in Scottish NSND revenue induced by a marginal, permanent proportional change in the level of earnings. No Scotland-specific figure is published, but the object is a standard one. It is what the public finance literature calls the \emph{revenue elasticity} or \emph{tax-to-base elasticity} of a progressive income tax, and it is the formal expression of fiscal drag: a schedule with nominal, unindexed thresholds mechanically raises revenue faster than the tax base whenever earnings grow, because a rising share of income is taxed at higher marginal rates. \citet{creedy2004} set out the UK approach, building on \citet{creedy2002} and \citet{johnson1989}, and it is a structural property of the schedule at a point in time rather than a parameter to be estimated from a time series. If every taxpayer's income rises by the fraction $\delta$ and the thresholds do not move, taxpayer $i$'s liability rises by $m_i y_i \delta$, where $m_i$ is the marginal rate at $y_i$, so
\begin{equation}\label{eq:elasticity}
\varepsilon \;=\;
\frac{\sum_i \omega_i\, m_i\, y_i}{\sum_i \omega_i\, T_i}
\;=\;
\frac{\text{income-weighted average marginal rate}}
     {\text{average rate on total income}},
\end{equation}
with $\omega_i$ the APS income weight, $T_i$ the liability under the prevailing Scottish schedule and the personal allowance taper carried explicitly. No assumption about threshold indexation is required, because \cref{eq:elasticity} is evaluated at the schedule as it stands. Progressivity alone puts $\varepsilon > 1$; the frozen higher, advanced and top rate thresholds are what keep it there as earnings grow \citep{sgreadyreckoner2026}.

On the pooled APS the income-weighted average marginal rate is 0.296 against an average rate on total income of 0.154, so $\varepsilon = 1.922$. Progressivity alone therefore very nearly doubles the revenue consequence of a given earnings differential.

\citet{garcia2026} compute the same elasticity by microsimulation for 21 European personal income tax systems and find it clusters between 1.7 and 2.0 in fourteen of them, with an unweighted average close to 1.8; \citet{price2015} report an OECD average of about 1.85 using representative-household calculations. $\varepsilon = 1.922$ sits inside that range, close to its upper end, which is what a six-band schedule with a tapered personal allowance and three widely-spaced higher bands should be expected to produce: \citet{garcia2026} find that elasticities are driven jointly by bracket progressivity and by the phase-out of allowances and credits, and the Scottish personal allowance taper is exactly such a phase-out, levied at an effective marginal rate one and a half times the headline band rate. \citet{immervoll2005} found comparatively low bracket-creep effects for the UK in a study using late-1990s data, but explains this by the contemporary UK schedule's few, wide bands. However, that structure is not the current Scottish one, which has added three further bands since 2018 specifically to increase progressivity at the top of the distribution, so a higher elasticity than Immervoll's UK estimate is the expected consequence of the reform rather than an anomaly. I am not aware of a published UK or Scottish estimate of this elasticity to benchmark against directly. The SFC's own income tax forecast instead grows a taxpayer-level micro-simulation of the Survey of Personal Incomes by projected earnings and re-applies the current schedule taxpayer by taxpayer, which captures fiscal drag endogenously and has no need of a scalar summary of it. \Cref{eq:elasticity} is a reduced-form statement of what that exercise does in aggregate, and the absence of a published SFC elasticity reflects that choice of method rather than the absence of the underlying effect.

Two departures from \cref{eq:elasticity} matter, and only one of them moves the number. The first is coverage. The APS records gross pay for employees only, so the computed distribution excludes self-employment and pension income and understates the upper tail, where marginal rates are highest. Grossed to a single wave, the implied aggregate liability is roughly three-fifths of the HMRC outturn, and \cref{tab:elasticity} reports the ratio so that $\varepsilon$ can be read against it. The direction of the resulting bias is worth being precise about. A fatter upper tail raises the numerator and the denominator of \cref{eq:elasticity} together, and since the marginal rate is bounded above by the top rate while the average rate is not, it raises the denominator proportionally more. The employee-only figure is therefore a mild overstatement, and I read 1.922 as the top of a narrow range rather than a point estimate. Because \cref{eq:elasticity} is a ratio of two identically weighted sums, it is invariant to the scale of the weights, and nothing in it turns on the grossing factor that the coverage row exposes.

The second is incidence. The flat pass-through of \cref{sec:fiscal} spreads the earnings differential evenly, whereas the mechanism of \cref{eq:tfp} concentrates it in exposed occupations, which are better paid and face higher marginal rates. Letting worker $i$'s earnings change in proportion to the exposure of their occupation, and normalising so the aggregate change is unchanged, gives the incidence-weighted elasticity
\begin{equation}\label{eq:elasticityE}
\varepsilon^{E} \;=\;
\frac{\sum_i \omega_i\, m_i\, y_i\, E_{j(i)}}
     {\bar{E}_y \sum_i \omega_i\, T_i},
\qquad
\bar{E}_y = \frac{\sum_i \omega_i\, y_i\, E_{j(i)}}{\sum_i \omega_i\, y_i},
\end{equation}
which exceeds $\varepsilon$ whenever exposure and the marginal rate both rise with earnings, as they do here. The correction is small: $\varepsilon^{E} = 1.974$ against $\varepsilon = 1.922$, so concentrating the earnings differential on exposed workers rather than spreading it evenly moves the fiscal figure by 2.7\%. \Cref{sec:fiscal} therefore carries the flat $\varepsilon$. 

\begin{table}[h]
\centering
\caption{Elasticity of Scottish NSND revenue to earnings}
\label{tab:elasticity}
\small
% Auto-generated by the analysis pipeline -- do not edit by hand.
% \input{} inside a table environment; requires \usepackage{booktabs}.
\begin{tabular}{l l}
\toprule
Quantity & Value \\
\midrule
Income-weighted average marginal rate & 0.296 \\
Average rate on total income & 0.154 \\
Elasticity of NSND revenue to earnings, eps & 1.922 \\
Incidence-weighted elasticity, eps\_E & 1.974 \\
Implied aggregate liability (GBP m) & 45786 \\
HMRC outturn, 2024-25 (GBP m) & 18635 \\
Coverage ratio (implied / outturn) & 2.46 \\
Scottish observations & 18,550 \\
Tax year of schedule & 2026-27 \\
\bottomrule
\end{tabular}

\begin{minipage}{0.92\textwidth}
\vspace{0.4em}\footnotesize
\textit{Notes:} Computed by \texttt{19\_nsnd\_elasticity.R} on the pooled APS
2022--2025 Scottish employee earnings distribution, weighted by the APS income weight, under the Scottish schedule of the stated tax year with the personal allowance taper carried explicitly. $\varepsilon$ is \cref{eq:elasticity} and $\varepsilon^{E}$ is \cref{eq:elasticityE}. The coverage row compares the aggregate liability implied by the weighted APS distribution with the HMRC outturn, and the elasticity should be read against that ratio.
\end{minipage}
\end{table}